\documentclass[aspectratio=169]{beamer}

% Use Sorbonne theme
\usepackage{sorbonne_theme}

% Title page information
\title{Econometrics Part 2}
\subtitle{The linear model in practice}
\author{Jose Angel Garcia Sanchez}
\institute{Sorbonne Data Analytics}
\date{} % left empty on purpose: beamer would otherwise print today's date

% Set the logo
\sorbonnelogo{imagen.png}

\begin{document}

% Title page (without logo in corner)
{
\setbeamertemplate{logo}{}
\begin{frame}
\titlepage
\end{frame}
}

% Table of contents
\begin{frame}{Outline}
\tableofcontents
\end{frame}

% Section 1: Multicollinearity
\section{Multicollinearity}

\begin{frame}{Multicollinearity}
\begin{itemize}
    \item This problem arises if some variable is equal to an exact linear combination of some other variables (e.g.\ if we have variables like \textit{income}, \textit{income after tax} and \textit{tax})
    \item This variable is unnecessary because it contains redundant information
    \item Worse, it prevents the computation of the OLS estimator, since $X$ is no longer full rank and we cannot compute $\hat{b} = (X'X)^{-1}X'y$
    \item Some software (e.g.\ Stata) detects it and arbitrarily removes one redundant variable
    \item statsmodels does \textbf{not}: it uses a pseudo-inverse and only prints a warning in the summary (``the design matrix is singular''), so always read it
\end{itemize}
\end{frame}

\begin{frame}{Near multicollinearity}
\begin{itemize}
    \item This problem arises if some variable is \textit{almost} equal to an exact linear combination of some other variables
    \item In other words, it means that they are very strongly correlated
    \item We \textit{can} compute $\hat{b} = (X'X)^{-1}X'y$
    \item But: $\hat{b}$ is very unstable and the estimator is not very reliable (large variance)
\end{itemize}
\end{frame}

\begin{frame}{How to detect near multicollinearity (1)}
\begin{itemize}
    \item The VIF (Variance Inflation Factor) indicates by how much the variance of a parameter $\hat{b}_j$ is inflated due to the correlation of variable $X_j$ with all the other explanatory variables
    \item Say our model has $k$ parameters: constant, $X_1$, $X_2$, \dots, $X_{k-1}$
    \item We run the regression of each $X_j$ on all the other explanatory variables, and save its $R^2$, called $R_j^2$
    \item Example: $X_1 = a_1 + a_2X_2 + \dots + a_{k-1}X_{k-1} + \varepsilon$
    \item Next, the VIF of each variable $j$ is computed as:
    \[
    \text{VIF}_j = \frac{1}{1-R_j^2}
    \]
\end{itemize}
\end{frame}

\begin{frame}{How to detect near multicollinearity (2)}
\begin{itemize}
    \item One can prove that in the original regression $y = b_0 + b_1X_1 + \dots + b_{k-1}X_{k-1} + u$:
    \[
    V(\hat{b}_j) = \frac{\sigma^2}{(N-1)\,\widehat{Var}(X_j)} \, \text{VIF}_j
    \]
    \item $\text{VIF}_j > 5$ (some use 10) indicates a high risk of multicollinearity between variable $j$ and the other variables
    \item Python: \texttt{from statsmodels.stats.outliers\_influence import variance\_inflation\_factor}, then \texttt{variance\_inflation\_factor(X.values, j)}
    \item Careful: $X$ must contain the constant column (\texttt{sm.add\_constant}), otherwise the VIFs are wrong
    \item If you don't see any big issue, do not remove variables: risk of omitted variable bias (see below)
\end{itemize}
\end{frame}

% Section 2: Handling categories
\section{Handling categories}

\begin{frame}{Dummy variables: 2 categories}
\begin{itemize}
    \item Let's say we have $N$ individual observations providing income, gender and the number of years of education
    \item We want to explain individual income $y$ by education $x$ and gender $z$: $y_i = a + b x_i + c z_i + u_i$
    \item $z$ is a \textit{categorical variable}, ideally coded with 0/1
    \item Say $z = 1$ codes for males
    \item $c$ is the average income difference between males and females with the same level of education
    \item It would be irrelevant to have both a male \textit{and} a female dummy: since male + female = 1 = the constant, it would cause \textit{exact multicollinearity} (the ``dummy variable trap'')
    \item Remark: having a categorical variable as the dependent variable requires other models, such as Probit and Logit (see later chapters)
\end{itemize}
\end{frame}

\begin{frame}{Dummy variables: more than 2 categories}
\begin{itemize}
    \item We want to explain life satisfaction $y$ by income $x$ and type of environment $z$: \textit{big city}, \textit{small city} or \textit{rural} (coded 1, 2 and 3)
    \item $z$ is still a \textit{categorical variable}, but with 3 categories: we choose a reference category, with respect to which parameters are computed, say \textit{rural}
    \item Model: $y_i = a + b x_i + c_1 z_{1i} + c_2 z_{2i} + u_i$, where $z_1$ and $z_2$ are dummies for \textit{big city} and \textit{small city}
    \item Do not keep the original coding (1, 2, 3), which would be meaningless in a regression
    \item pandas: \texttt{pd.get\_dummies(df["environment"], dtype=int)}, then drop the \textit{rural} column (\texttt{drop\_first=True} would drop \textit{big city}, the first category)
    \item statsmodels formula: \texttt{C(environment, Treatment(reference=3))}
\end{itemize}
\end{frame}

\begin{frame}{Dummy variables: outliers}
\begin{itemize}
    \item Let's say we have $N$ yearly observations for a country that experienced a shock
    \item We want to explain life expectancy $y$ by the number of doctors per capita $x$: $y_t = a + b x_t + u_t$
    \item Let's say that the shock took place in year $t_0$: this unusual observation is likely to bias the whole regression, but we don't want to delete it
    \item We will take this event into account with a dummy variable:
    \item $D_t = 1$ if $t = t_0$, $D_t = 0$ otherwise
    \item The model is now written: $y_t = a + b x_t + c D_t + u_t$
    \item $c$ represents the impact of the shock
    \item Note: the residual of year $t_0$ becomes 0, so $a$ and $b$ are estimated as if year $t_0$ were dropped
\end{itemize}
\end{frame}

% Section 3: Omitted variables
\section{Omitted variables}

\begin{frame}{Omitted variable bias}
\begin{itemize}
    \item Estimating the previous model with and without dummy $D$ can lead to changes even in parameters $a$ and $b$
    \item Omitting a relevant variable (here $D$) that is correlated with the included variables leads to \textit{omitted variable bias}, which affects all estimated parameters
    \item Proof: in the finite sample case with the Frisch-Waugh theorem (see Dormont), and proof involving consistency in later chapters (on endogeneity)
    \item This is a very important issue in applied work, so try not to forget important variables in models
\end{itemize}
\end{frame}

\begin{frame}{The Frisch-Waugh theorem}
\begin{itemize}
    \item Let $Z$ be a matrix of variables
    \item $M_Z = I - Z(Z'Z)^{-1}Z'$ is the orthogonal projection matrix on $L^{\perp}(Z)$
    \item Let's define 2 models:
    \begin{enumerate}
        \item $y = Xb_{(i)} + Zc + u$
        \item $M_Zy = M_ZXb_{(ii)} + M_Zu$
    \end{enumerate}
    \item The Frisch-Waugh theorem states that $\hat{b}_{(i)} = \hat{b}_{(ii)}$
    \item Important consequence: estimating the model with both sets of variables $X$ and $Z$ \textbf{or} with $X$ alone will in general not provide the same results
    \item This is called the ``omitted variable bias''
\end{itemize}

\textit{Under which condition will the results be the same?}
\end{frame}

\begin{frame}{Answer}
\begin{itemize}
    \item Estimating the model with both sets of variables $X$ and $Z$ \textbf{or} with $X$ alone will yield the same results only if $X$ and $Z$ are orthogonal
    \item i.e.\ $X'Z = 0$ (for centered variables: $X$ and $Z$ are totally uncorrelated)
    \item This rarely happens with observational data (it is the idea behind randomized experiments)
\end{itemize}
\end{frame}

\begin{frame}{A side effect of the FW theorem: centered variables}
\begin{itemize}
    \item Estimating the simple model $y_i = a + b x_i + u_i$ with or without centering observations (i.e.\ replacing each value $y_i$ by $y_i - \bar{y}$, same for $x$) will give the same result for $\hat{b}$
    \item Why: apply the FW theorem with $Z$ = the constant, since $M_Z$ simply centers the variables
    \item Centering observations will make the constant unnecessary
\end{itemize}
\end{frame}

% Section 4: Predictions in logs
\section{Predictions in logs}

\begin{frame}{Advantages of using logs}
\begin{enumerate}
    \item Taking the \textit{log} of the dependent variable can make the residuals look more ``normal''
    \item It can also reduce heteroskedasticity
    \item Parameters can be interpreted as percentages (semi-elasticities) or, if $x$ is also in logs, as elasticities
\end{enumerate}
\end{frame}

\begin{frame}{About the log transformation}
\begin{itemize}
    \item The average predicted $y$ should be equal to the average observed $y$
    \item Here, we can compute a prediction of $\log(y)$
    \item Can we get back to the predicted $y$ (in levels) simply by taking the exponential of the predicted $\log(y)$?
\end{itemize}
\end{frame}

\begin{frame}{Log and predictions (1)}
\begin{itemize}
    \item Let's use model $\forall i$, $\log(y_i) = a + bx_i + u_i$ (what matters is that we take the \textit{log} of $y$)
    \item Taking the exponential:\\
    $y_i = \exp(a + bx_i + u_i) = \exp(a + bx_i)\exp(u_i)$
    \item So if the $u$'s are independent from the $x$'s, we get:\\
    $E[y_i] = E[\exp(a + bx_i)\exp(u_i)] = E[\exp(a + bx_i)]\,E[\exp(u_i)]$
    \item We usually consider the $x$'s as fixed, so:\\
    $E[y_i] = \exp(a + bx_i)\,E[\exp(u_i)]$
    \item By Jensen's inequality, $E[\exp(u_i)] > \exp(E[u_i]) = 1$: simply taking the exponential underestimates $y$
    \item A prediction for $y$ is thus: $\hat{y}_i = \exp(\hat{a} + \hat{b}x_i)\,\hat{E}[\exp(u_i)]$
    \item But how to estimate $E[\exp(u_i)]$ (the retransformation parameter)?
\end{itemize}
\end{frame}

\begin{frame}{Log and predictions (2)}
It can be proven that if $u \hookrightarrow N(0,\sigma^2)$, then $E[\exp(u_i)] = \exp(\frac{\sigma^2}{2})$. This can be estimated using:
\[
\hat{\Psi} = \exp\left(\frac{\hat{\sigma}^2}{2}\right)
\qquad \text{so that} \qquad
\hat{y}_i = \exp(\hat{a} + \hat{b}x_i)\,\hat{\Psi}
\]

To get back to levels when using a prediction in \textit{logs}, we cannot simply take the exponential of the prediction: there is an additional retransformation parameter to take into account.

The preceding formula holds only if the $u$'s are normal: if their distribution is unknown but homoskedastic (same variance), we can use Duan's non-parametric\footnote{Because it makes no assumption on the type of distribution} \textit{smearing} estimator:
\[
\hat{\Psi} = \frac{1}{N}\sum_{i=1}^{N} \exp(\hat{u}_i)
\]
\end{frame}

% Section 5: Influential observations
\section{Influential observations}

\begin{frame}{Are all observations given the same weight?}
We have:
\[\hat{b}_1 = \frac{\sum(x_i - \bar{x})(y_i - \bar{y})}{\sum(x_i - \bar{x})^2}\]

So:
\[\hat{b}_1 = \frac{\sum(x_i - \bar{x})^2 \frac{(y_i-\bar{y})}{(x_i-\bar{x})}}{\sum(x_i - \bar{x})^2} = \sum p_i\frac{(y_i - \bar{y})}{(x_i - \bar{x})}\]

Calling $p_i = \frac{(x_i - \bar{x})^2}{\sum(x_i - \bar{x})^2}$.

$\frac{(y_i-\bar{y})}{(x_i-\bar{x})}$ is the slope of the line drawn from the point corresponding to individual $i$ to the sample average, and $p_i$ is an increasing function of $|x_i - \bar{x}|$.

Estimate $\hat{b}_1$ is thus highly influenced by points with extreme $x$ values (see Anscombe's quartet: four datasets with the same regression line).
\end{frame}

\begin{frame}{Outliers (1)}
\begin{itemize}
    \item Outlier: observation with a large residual compared to other observations
    \item They may indicate a poor goodness of fit
    \item How to detect them: using a residuals plot
    \item Goal: estimate how far the observed value is from the predicted value
    \item Or more rigorously: using the \textit{studentized} residuals, which adjust for the \textit{leverage} of that particular point
\end{itemize}
\end{frame}

\begin{frame}{Outliers (2)}
\begin{itemize}
    \item Outliers are observations that have a large residual
    \item However, if the observation has a high leverage, it has pulled the predicted value towards itself
    \item In other words, the variances of the residuals differ, even though the variances of the true errors are all equal to each other
    \item Consequence: need for studentization
    \item The \textit{studentized} residuals adjust for the leverage of that particular point (see what $h$ is below)
\end{itemize}

\[rstud_i = \frac{\hat{u}_i}{\hat{\sigma}\sqrt{1 - h_i}}\]

This is the \textit{internally} studentized residual; the \textit{externally} studentized one uses $\hat{\sigma}_{(-i)}$, estimated without observation $i$.
\end{frame}

\begin{frame}{Leverage (1)}
\begin{itemize}
    \item High leverage observation: has an explanatory variable that takes a value very different from the sample mean
    \item In a simple regression model, it could be (for example) at the far right and would influence the slope of the regression line
    \item How to detect them: the ``hat statistic''
\end{itemize}
\end{frame}

\begin{frame}{Leverage (2)}
\begin{itemize}
    \item We know that $\hat{y} = P_Xy$
    \item $P_X$ can be called the ``hat'' matrix (it puts a hat on $y$)
    \item Let's call $h_i$ the $i^{th}$ element of the diagonal of the hat matrix
    \item $h_i$ is the \textit{leverage} of observation $i$: if it is large, then observation $i$ has a high leverage
    \item $0 \leq h_i \leq 1$ (and $h_i \geq 1/N$ when the model has a constant)
    \item $\sum_i h_i = k$, so the average leverage is $\bar{h} = k/N$
\end{itemize}
\end{frame}

\begin{frame}{Leverage (3)}
\begin{itemize}
    \item We know that $\hat{u} = (I - P_X)u$, so $V(\hat{u}) = \sigma^2(I - P_X)$
    \item So: $V(\hat{u}_i) = \sigma^2(1 - h_i)$
    \item The greater the leverage $h_i$, the smaller the variance of the residual for observation $i$
    \item Meaning: high leverage observations tend to bring the regression line close to them
    \item High leverage observations are not necessarily influential, but low leverage observations are rarely very influential
    \item For the simple linear model: $h_i = \frac{1}{N} + \frac{(x_i-\bar{x})^2}{\sum(x_i-\bar{x})^2}$
\end{itemize}
\end{frame}

\begin{frame}{Influence measures (1)}
\begin{itemize}
    \item Influential observation: if removed from the regression, estimated parameters vary a lot
    \item How to detect them: the \textit{Dfbeta} statistic
    \item For every variable $k$, its \textit{Dfbeta} measures how much its coefficient would change if we removed one observation
    \item It is computed for every observation $i$ and every coefficient $k$ (an $N \times k$ table)
    \item The change in $\hat{b}_k$ if observation $i$ is dropped is:
\end{itemize}

\[\text{Dfbeta}_{i,k} = \frac{\hat{b}_k - \hat{b}_{k,(-i)}}{\hat{\sigma}_{k,(-i)}}\]

with $\hat{b}_{k,(-i)}$ and $\hat{\sigma}_{k,(-i)}$ computed without observation $i$. It represents the number of standard errors by which $\hat{b}_k$ changes when observation $i$ is removed.
\end{frame}

\begin{frame}{Influence measures (2)}
\begin{itemize}
    \item \textit{Dffit}: measures how much the prediction for observation $i$ changes when the estimation is made \textit{without} observation $i$
    \item $\text{Dffit}_i = \frac{\hat{y}_i-\hat{y}_{i,(-i)}}{\hat{\sigma}_{(-i)}\sqrt{h_i}}$
    \item \textit{Covratio} statistic: measures the change in the determinant of the covariance matrix of the estimates when deleting the $i^{th}$ observation (indicates loss in precision)
    \item \textit{Cook's distance}: measures the influence of observation $i$ on the model as a whole, not on a single coefficient (it combines all the \textit{Dfbeta}'s)
    \item $D_i = \frac{(\hat{b}-\hat{b}_{(-i)})'(X'X)(\hat{b}-\hat{b}_{(-i)})}{k\hat{\sigma}^2}$
\end{itemize}
\end{frame}

\begin{frame}{Cutoffs}
\begin{itemize}
    \item Leverage $h_i$ of observation $i$: the average is $\bar{h} = k/N$. If $h_i > 2k/N$ (or $3k/N$), observation $i$ has a high leverage and might be influential
    \item Change in predicted value \textit{Dffit}: observation $i$ is influential if $|\text{Dffit}_i| > 2$ (general cutoff) or $> 2\sqrt{k/N}$
    \item Change in a parameter \textit{Dfbeta}: observation $i$ is influential if $|\text{Dfbeta}_{i,k}| > 2$ (general cutoff) or $> 2/\sqrt{N}$ (adj.\ for large sample)
    \item Impact on the whole model (Cook's distance $D$): observation $i$ is influential if $D_i > 1$ (some use $4/N$)
    \item \textit{Covratio} statistic: observation $i$ is influential if $|\text{covratio}_i - 1| > 3k/N$
    \item Python: \texttt{infl = res.get\_influence()}; \texttt{infl.summary\_frame()} gives leverage, Dfbeta, Dffit and Cook's distance for each observation, and \texttt{infl.cov\_ratio} the covratio
\end{itemize}
\end{frame}

\end{document}
